% !TEX program = lualatex
\documentclass[professionalfonts]{beamer}
\newcommand{\slidemode}{1}  % カラー投影モード
\usepackage{beamer_template}
\usepackage{enumerate}

\newcommand{\LectureNum}{3}
\newcommand{\LectureTitle}{微分法則と積分法則}
\newcommand{\CourseName}{応用数学}
\renewcommand{\TermName}{前期}

\begin{document}

\begin{frame}{第3回　微分法則と積分法則}
  \begin{block}{今回の内容}
  微分法則と積分法則
  \end{block}
  \vfill\hfill{\scriptsize \textcolor{gray}{第2章 ラプラス変換　§1}}
\end{frame}

\begin{frame}{定理：原関数の微分法則}
  \begin{block}{}
  \begin{align*}
    L[f'(t)]    &= sF(s) - f(0) \\
    L[f''(t)]   &= s^{2}F(s) - sf(0) - f'(0) \\
    L[f^{(n)}(t)] &= s^n F(s) - s^{n-1}f(0) - s^{n-2}f'(0) - \ldots - f^{(n-1)}(0)
  \end{align*}
  \end{block}
\end{frame}

\begin{frame}{定理：像関数の高次微分}
  \begin{block}{}
  \[
    L[t^n f(t)] = (-1)^n F^{(n)}(s)
  \]
  \end{block}
\end{frame}

\begin{frame}{例題9}
  \begin{exampleblock}{問題}
    像関数の高次微分を用いて次の問いに答えよ\\[2pt]
    \begin{description}
      \item[(1)] $L[te^{\alpha t}]$ を求めよ（$\alpha$ は定数）
      \item[(2)] $L[t\sin(\omega t)]$ を求めよ（$\omega$ は定数）
    \end{description}
  \end{exampleblock}
\end{frame}

\begin{frame}{例題9(1)　解答}
  $f(t) = e^{\alpha t}$ とおくと $F(s) = L[e^{\alpha t}] = 1/(s-\alpha)$\\[4pt]
  像関数の高次微分 $L[t^n f(t)] = (-1)^n F^{(n)}(s)$ を $n=1$ で適用\\[4pt]
  $F'(s) = d/ds[1/(s-\alpha)] = -1/(s-\alpha)^{2}$\\[4pt]
  $L[te^{\alpha t}] = (-1)^{1}\cdot F'(s) = 1/(s-\alpha)^{2}$
  \\[8pt]\textcolor{red}{\textbf{答}}：$L[te^{\alpha t}] = 1/(s-\alpha)^{2} \ (s > \alpha)$
\end{frame}

\begin{frame}{例題9(2)　解答}
  $f(t) = \sin(\omega t)$ とおくと $F(s) = L[\sin(\omega t)] = \omega/(s^{2}+\omega^{2})$\\[4pt]
  像関数の高次微分を $n=1$ で適用\\[4pt]
  $F'(s) = \omega\cdot\dfrac{-(s^{2}+\omega^{2})'}{(s^{2}+\omega^{2})^{2}} = \omega\cdot\dfrac{-2s}{(s^{2}+\omega^{2})^{2}} = \dfrac{-2\omega s}{(s^{2}+\omega^{2})^{2}}$\\[4pt]
  $L[t\sin(\omega t)] = (-1)^{1}\cdot F'(s) = 2\omega s/(s^{2}+\omega^{2})^{2}$
  \\[8pt]\textcolor{red}{\textbf{答}}：$L[t\sin(\omega t)] = 2\omega s/(s^{2}+\omega^{2})^{2} \ (s > 0)$
\end{frame}

\begin{frame}{演習問題}
  \begin{exampleblock}{自分で解いてみよう}
  \begin{description}
    \item[問13] $L[t \cos(\omega t)]$ を求めよ
    \item[問14] $f'(t)+2f(t)=t, f(0)=0$ を満たすとき、 $L[f(t)]$ を求めよ
  \end{description}
  \end{exampleblock}
\end{frame}

\begin{frame}{問13　解答}
  像関数の高次微分 $: L[t^n f(t)] = (-1)^n F^{(n)}(s)$ を $n=1$ で使用\\[4pt]
  $F(s) = L[\cos(\omega t)] = s/(s^{2}+\omega^{2})$\\[4pt]
  $F'(s) = d/ds [s/(s^{2}+\omega^{2})] = (s^{2}+\omega^{2}-2s^{2})/(s^{2}+\omega^{2})^{2} = (\omega^{2}-s^{2})/(s^{2}+\omega^{2})^{2}$\\[4pt]
  $L[t \cos(\omega t)] = (-1)^{1}\cdot F'(s) = -F'(s)$
  \\[8pt]\textcolor{red}{\textbf{答}}：$L[t \cos(\omega t)] = (s^{2}-\omega^{2})/(s^{2}+\omega^{2})^{2}$
\end{frame}

\begin{frame}{問14　解答}
  両辺をラプラス変換 $: (sF(s)-f(0)) + 2F(s) = 1/s^{2}$\\[4pt]
  $f(0)=0$ なので $(s+2)F(s) = 1/s^{2}$
  \\[8pt]\textcolor{red}{\textbf{答}}：$L[f(t)] = F(s) = 1/(s^{2}(s+2))$
\end{frame}

\begin{frame}{例題10}
  \begin{exampleblock}{問題}
    $L[t^n e^t]$ を求めよ（ n は正の整数）\\[2pt]
  \end{exampleblock}
\end{frame}

\begin{frame}{例題10　解答}
  $f(t) = e^{t}$ とおくと $F(s) = L[e^{t}] = 1/(s-1)$\\[4pt]
  像関数の高次微分 $L[t^n f(t)] = (-1)^n F^{(n)}(s)$ を適用\\[4pt]
  $F^{(n)}(s) = d^n/ds^n[1/(s-1)] = (-1)^n n!/(s-1)^{n+1}$\\[4pt]
  $L[t^n e^{t}] = (-1)^n\cdot(-1)^n n!/(s-1)^{n+1}$
  \\[8pt]\textcolor{red}{\textbf{答}}：$L[t^n e^{t}] = n!/(s-1)^{n+1}  (s > 1)$
\end{frame}

\begin{frame}{演習問題}
  \begin{exampleblock}{自分で解いてみよう}
  \begin{description}
    \item[問15] $L[t^n]$ を求めよ（ n は正の整数）（像関数の高次微分を使用）
  \end{description}
  \end{exampleblock}
\end{frame}

\begin{frame}{問15　解答}
  $f(t) = 1$ とおくと $F(s) = L[1] = 1/s$\\[4pt]
  $F^{(n)}(s) = d^n/ds^n [1/s] = (-1)^n n!/s^{n+1}$\\[4pt]
  $L[t^n\cdot1] = (-1)^n F^{(n)}(s) = (-1)^n\cdot(-1)^n n!/s^{n+1}$
  \\[8pt]\textcolor{red}{\textbf{答}}：$L[t^n] = n!/s^{n+1}  (s > 0)$
\end{frame}

\begin{frame}{定理：原関数の積分法則}
  \begin{block}{}
  \[
    L[\int_{0}^t f(\tau)d\tau] = F(s)/s
  \]
  \end{block}
\end{frame}

\begin{frame}{定理：像関数の積分}
  \begin{block}{}
  \[
    L[f(t)/t] = \int_{s}^\infty F(u)du
  \]
  \end{block}
\end{frame}

\begin{frame}{例題11}
  \begin{exampleblock}{問題}
    $L[\sin(\omega t)/t]$ を求めよ（ $\omega$ は正の定数）\\[2pt]
  \end{exampleblock}
\end{frame}

\begin{frame}{例題11　解答}
  像関数の積分 $: L[f(t)/t] = \int_{s}^\infty F(u)du$ を利用\\[4pt]
  $L[\sin(\omega t)] = \omega/(u^{2}+\omega^{2})$ なので\\[4pt]
  $L[\sin(\omega t)/t] = \int_{s}^\infty \omega/(u^{2}+\omega^{2}) du = [\arctan(u/\omega)]_{s}^\infty$\\[4pt]
  $= \pi/2 - \arctan(s/\omega)$
  \\[8pt]\textcolor{red}{\textbf{答}}：$L[\sin(\omega t)/t] = \pi/2 - \arctan(s/\omega) = \arctan(\omega/s) \ (s > 0)$
\end{frame}

\begin{frame}{演習問題}
  \begin{exampleblock}{自分で解いてみよう}
  \begin{description}
    \item[問16] $L[(e^{2t} - e^t)/t]$ を求めよ
  \end{description}
  \end{exampleblock}
\end{frame}

\begin{frame}{問16　解答}
  像関数の積分 $: L[f(t)/t] = \int_{s}^\infty F(u)du$ を利用\\[4pt]
  $f(t) = e^{2t} - e^t$ とおくと $F(u) = L[e^{2t}-e^t]|_{s=u} = 1/(u-2) - 1/(u-1)$\\[4pt]
  $L[(e^{2t}-e^t)/t] = \int_{s}^\infty (1/(u-2) - 1/(u-1)) du$\\[4pt]
  $= [\ln(u-2) - \ln(u-1)]_{s}^\infty$\\[4pt]
  $= \lim_{u \to \infty} \ln((u-2)/(u-1)) - \ln((s-2)/(s-1))$\\[4pt]
  $= 0 - \ln((s-2)/(s-1)) = \ln((s-1)/(s-2))$
  \\[8pt]\textcolor{red}{\textbf{答}}：$L[(e^{2t}-e^t)/t] = \ln((s-1)/(s-2)) \ (s > 2)$
\end{frame}

\end{document}
